\documentclass[../../../include/open-logic-section]{subfiles}

\begin{document}

\olfileid{sth}{choice}{woproblem}
\olsection{حُسنِ ترتیب دا مسئلہ}

صاف اے کہ ساڈے حُسنِ ترتیب نوں منّن یا نہ منّن اُتے کافی کجھ ٹکدا اے۔ پر ایس اصول دی بحث دا دھیان اک برابر دے اصول، چُون دے مسلّمے، ول رہیا اے۔ سو ہُن اسیں اوہدی گل کراں گے (تے دوہاں دی برابری ثابت کراں گے)۔

\citeyear{Cantor1883} وچ کانتور نے حُسنِ ترتیب دے مسلّمے دی تائید کیتی تے اینوں «سوچ دا اک قانون جیہڑا مینوں بنیادی، نتیجیاں وچ امیر، تے ہر تھاں لاگو ہون کرکے خاص دھیان جوگا لگدا اے» آکھیا (\citeauthor{Potter2004} \citeyear[p.~243]{Potter2004} وچ ایہ قول دِتا اے)۔ پر آخر کانتور نوں یقین ہو گیا کہ ایس «مسلّمے» دا ثبوت چاہیدا اے۔ ریاضی دان وی ایس نتیجے اُتے پہنچے۔

زرمَیلو نے \citeyear{Zermelo1904} وچ ایہ مسئلہ «حل» کیتا۔ اوہدا حل سمجھاؤن لئی سانوں کجھ تعریفاں چاہیدیاں نیں۔

\begin{defn}
اک فنکشن $f$ اک \emph{چُون والا فنکشن} اے اُدوں تے صرف اُدوں جدوں $f(x) \in x$ سارے $x \in \dom{f}$ لئی۔ اسیں آکھدے آں کہ $f$ اک \emph{$A$ لئی چُون والا فنکشن} اے اُدوں تے صرف اُدوں جدوں $f$ اک چُون والا فنکشن اے جیہدے لئی $\dom{f} = A \setminus \{\emptyset\}$۔
\end{defn}

سہج سمجھ مطابق، ہر (غیر خالی) سیٹ $x \in A$ لئی، $A$ دا چُون والا فنکشن اک خاص رکن، $f(x)$، $x$ وچوں \emph{چُندا} اے۔ سو چُون دا مسلّمہ ایہ اے:

\begin{axiom}[چُون]
	ہر سیٹ لئی اک چُون والا فنکشن اے۔
\end{axiom}

زرمَیلو نے وکھایا کہ چُون توں حُسنِ ترتیب لازم آؤندی اے، تے اُلٹے پاسے وی:

\begin{thm}[$\ZF$ وچ]\ollabel{thmwochoice}
حُسنِ ترتیب تے چُون دے مسلّمے برابر نیں۔
\end{thm}

\begin{proof}
\emph{کھبّے توں سجّے ول۔} $A$ سیٹاں دا اک سیٹ ہووے۔ تاں اتحاد دے مسلّمے نال $\bigcup A$ موجود اے، تے حُسنِ ترتیب دے مسلّمے نال کوئی $<$ اے جیہڑا $\bigcup A$ نوں حُسنِ ترتیب دیندا اے۔ ہُن $f(x) = \text{$<$ مطابق سبھ توں چھوٹا رکن، جیہڑا ایس سیٹ وچ اے }x$ ہووے۔ ایہ $A$ لئی چُون والا فنکشن اے۔

\emph{سجّے توں کھبّے ول۔} $A$ مقرر کرو۔ چُون دے مسلّمے نال اک چُون والا فنکشن، $f$، اے، جیہڑا $\Pow{A} \setminus \{\emptyset\}$ لئی اے۔ لامتناہی درجیاں والی بازگشت ورت کے اک فنکشن ایویں تعریف کرو:
\begin{align*}
	g(0) &= f(A)\\
	g(\alpha) &=
		\begin{cases}
			\text{رُک جاؤ!{}} &\text{جے }A = \funimage{g}{\alpha}\\
			f(A \setminus \funimage{g}{\alpha}) & \text{ہور صورت وچ}\\
		\end{cases}
\end{align*}
«رُک جاؤ!» بس اوہ تعریف مختصر کرن لئی آکھیا اے جیہڑی ہور طرح لمّی ہو جانی سی۔ یعنی جدوں پہلی واری $A =
\funimage{g}{\alpha}$ ہووے، تاں $g(\delta) = A$ سارے $\delta \leq \alpha$ لئی ہووے۔ ہُن پہلاں اسیں صرف اینا یقین کر سکدے آں کہ ایس نال اک \emph{ٹرم} تعریف ہُندی اے (\olref[sth][spine][recursion]{transrecursionschema} توں بعد دیاں گلاں ویکھو)۔ پر اسیں وکھاواں گے کہ ساڈے کول واقعی اک فنکشن اے۔

کیوں جے $f$ اک چُون والا فنکشن اے، ہر $\alpha$ لئی (جدوں تعریف ہوئی ہووے) سانوں $g(\alpha) =
f(A \setminus \funimage{g}{\alpha}) \in A \setminus
\funimage{g}{\alpha}$ ملدا اے، یعنی $g(\alpha) \notin \funimage{g}{\alpha}$۔ سو جے $g(\alpha) = g(\beta)$، تاں $g(\beta) \notin
\funimage{g}{\alpha}$، یعنی $\beta \notin \alpha$، تے ایویں ای $\alpha \notin \beta$۔ سو تِن صورتاں دے اصول مطابق $\alpha = \beta$۔ سو $g$ !!{injective} اے۔

اگے دھیان دیو کہ اسیں واقعی رُکدے آں!{} یعنی اک (سبھ توں چھوٹا) ترتیبی عدد $\alpha$ اے جیہدے لئی $A = g[\alpha]$۔ کیوں جے منّ لو ایہ گل نہ ہووے، تاں $g$ دے !!{injective} ہون کرکے $\cardless{\alpha}{\Pow{A} \setminus \{\emptyset\}}$ ہر ترتیبی عدد $\alpha$ لئی ہووے گا، تے ایہ \olref[hartogs]{HartogsLemma} نال ٹکراؤندا اے۔ ایس توں $\ran{g} = A$ وی۔

ایہناں گلاں نوں جوڑیے تاں $g$ کسے ترتیبی عدد توں $A$ ول اک !!a{bijection} اے۔ ہُن $g$ نوں $A$ نوں حُسنِ ترتیب دین لئی ورت سکدے آں۔
\end{proof}

سو حُسنِ ترتیب تے چُون دے مسلّمے اکٹھے منّے یا ردّ کیتے جاندے نیں۔ پر سوال رہندا اے کہ اوہ منّے جان یا ردّ کیتے جان؟

\end{document}
